% SPDX-License-Identifier: 0BSD
\section{\texorpdfstring{Worked \(d\bar d\to Zgg\) Example}{Worked d dbar -> Z g g Example}}
\label{sec:zgg-example}

This example illustrates how a familiar tree-level process is represented by
reusable off-shell currents.  It also fixes the momentum, helicity, and colour
conventions used in the performance tables.

\subsection{External legs and conventions}

The physical process is
\begin{equation}
  d(p_1)\,\bar d(p_2) \longrightarrow Z(p_3)\,g(p_4)\,g(p_5),
  \qquad p_1+p_2=p_3+p_4+p_5 .
  \label{eq:zgg-physical-process}
\end{equation}
The current recursion uses the all-outgoing convention
\begin{equation}
  q_1=-p_1,\quad q_2=-p_2,\quad q_i=p_i\ (i=3,4,5),
  \qquad \sum_{i=1}^{5}q_i=0 .
  \label{eq:zgg-all-outgoing}
\end{equation}
External momenta are supplied in the physical convention of
Eq.~\eqref{eq:zgg-physical-process} and crossed to the all-outgoing convention
before the recursion is evaluated.

\begin{center}
\small
\begin{tabular}{@{}cL{0.90in}L{1.05in}L{1.00in}cL{1.35in}@{}}
\toprule
\textbf{label} & \textbf{physical leg} & \textbf{outgoing state} &
\textbf{momentum} & \textbf{states} & \textbf{colour role} \\
\midrule
$1$ & incoming $d$ & outgoing $\bar d$ & $q_1=-p_1$ & $2$ &
  $\overline{\mathbf 3}$, line end \\
$2$ & incoming $\bar d$ & outgoing $d$ & $q_2=-p_2$ & $2$ &
  $\mathbf 3$, line start \\
$3$ & outgoing $Z$ & outgoing $Z$ & $q_3=p_3$ & $3$ & singlet \\
$4$ & outgoing $g$ & outgoing $g$ & $q_4=p_4$ & $2$ & adjoint insertion \\
$5$ & outgoing $g$ & outgoing $g$ & $q_5=p_5$ & $2$ & adjoint insertion \\
\bottomrule
\end{tabular}
\end{center}

There are \(2\times2\times3\times2\times2=48\) physical helicity
configurations.  Twenty-four vanish identically because the massless quark line
permits only chirality-compatible initial-state helicities; they remain
explicit zero entries in resolved results.

At leading colour the \(Z\) is absent from the colour word.  Complete coverage
contains the two orderings
\begin{equation}
  c_0=(2,4,5,1),\qquad c_1=(2,5,4,1).
  \label{eq:zgg-physical-flows}
\end{equation}
Both are retained during generation, so either flow can be selected at
evaluation time.

\subsection{Reusable current recursion}

For a nonempty source subset \(I\), define its momentum by
\begin{equation}
  q_I=\sum_{i\in I}q_i .
  \label{eq:zgg-subset-mask}
\end{equation}
A one-leg current is the appropriate external spinor, polarization vector, or
scalar wavefunction.  Composite currents obey the binary recursion
\begin{equation}
  J_{t,\Omega}^{a}(I)
  =P_t{}^{a}{}_{b}(q_I)
   \sum_{\substack{I=A\mathbin{\dot\cup}B\\
                   A,B\ne\varnothing}}
   \sum_{v\,:\,t_A t_B\to t}
   \kappa_v\,
   V_v{}^{b}{}_{cd}(q_A,q_B)
   J_{t_A,\Omega_A}^{c}(A)
   J_{t_B,\Omega_B}^{d}(B).
  \label{eq:zgg-current-recursion}
\end{equation}
Here \(t\) is the particle state, \(P_t\) is its propagator, \(V_v\) is an
allowed interaction, and \(\Omega\) collects the quantum numbers required to
identify the current.  Only combinations consistent with the model,
perturbative order, and colour ordering enter the sum.

Representative terms along the quark line are
\begin{align}
 J_d(2,5) &={P}_d\,V_{dgd}\bigl[J_d(2),J_g(5)\bigr],
 \label{eq:zgg-j25}\\
 J_g(4,5) &={P}_g\,V_{ggg}\bigl[J_g(4),J_g(5)\bigr],
 \label{eq:zgg-j45}\\
 J_d(2,4,5) &={P}_d\Bigl(
     V_{dgd}\bigl[J_d(2,5),J_g(4)\bigr]
    +V_{dgd}\bigl[J_d(2,4),J_g(5)\bigr]
    +V_{dgd}\bigl[J_d(2),J_g(4,5)\bigr]
  \Bigr),
 \label{eq:zgg-j245}\\
 J_d(2,3,4,5) &={P}_d\!\sum_{A\dot\cup B=\{2,3,4,5\}}
   V\bigl[J(A),J(B)\bigr].
 \label{eq:zgg-j2345}
\end{align}
Equation~\eqref{eq:zgg-j2345} includes every admissible placement of the
\(Z\) emission and both gluon insertions.  Ordered currents such as
\(J_d(2,5,4)\) and \(J_d(2,4,5)\) remain distinct even though their momentum
subset is the same.

The final closure joins the endpoint current to source label \(1\):
\begin{equation}
  {\cal A}_{h,c}^{(r)}
  =C^{(r)}_{\dot a a}\,
   J_{\bar d}^{\dot a}(1)\,
   J_d^a(2,3,4,5),
  \qquad
  {\cal A}_{h,c}=\sum_{r\in G(h,c)}{\cal A}_{h,c}^{(r)} .
  \label{eq:zgg-amplitude-root}
\end{equation}
The set \(G(h,c)\) contains contributions that must be summed coherently.
Identical subcurrents are evaluated once and reused wherever their complete
quantum-number and ordering labels agree.  This sharing is the essential
difference between the current DAG and an independent evaluation of every
Feynman diagram.

\subsection{Resolved observable and normalization}

The two leading-colour execution layouts evaluate the same complete set of
helicity and flow components.  One is optimized for a selected flow with summed helicities;
the other shares work across all flows for a selected helicity.  Their resolved
components and their summed matrix element are numerically identical.

For the built-in Standard Model, the normalization used in this example is
\begin{equation}
  |{\cal M}|^2
  ={\cal N}\sum_{h,c}w_{h,c}\,
     \left|\sum_{r\in G(h,c)}{\cal A}_{h,c}^{(r)}\right|^2,
  \qquad
  {\cal N}=\frac{27\,(4\pi\alpha_s)^2(8\pi\alpha_{\rm EW})}
                  {36\times2}.
  \label{eq:zgg-normalization}
\end{equation}
The denominator averages over the incoming colour and spin states, while the
factor \(2=2!\) accounts for the identical final-state gluons.  Coupling
powers, colour weights, coherent groups, and symmetry factors are retained
separately so that changing numerical model parameters does not alter the
process definition.
